\documentclass[12pt]{article}
\usepackage[utf8]{inputenc}
\usepackage[T1]{fontenc}
\usepackage{hyperref}
\title{Ura Language Resources for K-12 Curriculum Development}
\author{Compiled by Claude Sonnet 5 (claude-sonnet-5) for Andr\'as Kornai}
\date{}
\begin{document}
\maketitle
\section{Introduction}
Ura (also Uram\"at; ISO 639-3 uro, Glottocode urap1240) is spoken in East New Britain Province, Papua New Guinea, by about 1,900 people (1991 count). Glottolog places it in the Baining group and lists it as shifting. Do not confuse it with the Ura language of Erromango (Vanuatu). Search results indexed some Lynch (1983) works under Erromango, so those are not cited here.
\section{Primary Educational Resources}
None found.
\section{Linguistic and Reference Materials}
\begin{itemize}
\item Ura language (Papua New Guinea) (Wikipedia): \url{https://en.wikipedia.org/wiki/Ura_language_(Papua_New_Guinea)}
\item Glottolog entry: \url{https://glottolog.org/resource/languoid/id/urap1240}
\item Endangered Languages Project entry: \url{https://endangeredlanguages.com/lang/3414}
\item Thesis \emph{Topics in Ura Phonology and Morphophonology, with Lexicographic Application} (University of Canterbury, 2007; rights reserved), Digital Pasifik record: \url{https://digitalpasifik.org/items/644428}
\end{itemize}
\section{Religious and Community Resources}
None found.
\section{Digital and Audio Resources}
\begin{itemize}
\item Wikitongues language profile: \url{https://wikitongues.org/languages/uro/}
\end{itemize}
\section{Curriculum Development Context}
Little is available. The language is not known to be taught in schools. The thesis record may point to lexicographic material, but the thesis may be inaccessible, and its language identification should be verified.
\end{document}